\section{Security Envelope: Details}
\label{app:envelope}

This appendix holds the analyses behind Section~\ref{sec:discussion}: the exposure before the lock at mainnet cadence, the choice of the down threshold, and the proving latency of the recovery proof.

\subsection*{Exposure before the lock at mainnet cadence}

Section~\ref{sec:discussion} converts the measured lock time to mainnet cadence by
scaling. To check that conversion we drove the real FSM with Bitcoin blocks and a model of
the checkpoint cycle of the anchor tracker: one checkpoint is pending at a time,
it is included in the next Bitcoin block, it counts as confirmed with
\texttt{KDeepFinality}+1 confirmations, the next one is submitted at once, and
the committed anchored height follows the tracker after a lag of two to three
Engram blocks, the one free parameter. At testbed cadence the model reproduces
what we measured: the checkpoint pause reaches \textsc{Sovereign} after 129\,s
(5th to 95th percentile 91 to 149\,s; measured 125\,s in RQ2 and $141.6 \pm
13.8$\,s in S2), and the quiet phase of the live sweeps is 0.80--0.84 in
\textsc{Anchored} at $k{=}2$ and 1.00 at $k{=}4$ (measured 0.79--0.87 and 1.00).
At mainnet cadence the same fault reaches \textsc{Sovereign} after 49.6 minutes
with steady 600\,s blocks (41 to 58) and 46.7 minutes with exponential block
intervals (18 to 79), which is about five Bitcoin blocks and 2{,}300--2{,}500
Engram blocks. This is lower than the 62 minutes and 3{,}100 blocks of
Table~\ref{tab:exposure}, which scales the measured 125\,s by thirty: the table is
the plain conversion and the model follows the checkpoint cycle, so the two
bracket the lock time at mainnet cadence.

The healthy gap is not flat. It follows the checkpoint cycle, 2, 3, 4, 5 and
back to 2 every three Bitcoin blocks, and in the first 18\,s of the five
composite runs, before any fault, it sat at 5, the \texttt{SuspiciousThreshold},
in 3--16\% of the samples, for about two Engram blocks each time. This has two
consequences. The threshold cannot simply be lowered, since the healthy cycle
already touches it. And part of the flapping at $k{=}2$ in the live sweeps is the
cycle, not the wide-area noise: the model gives one demotion per cycle, 1{,}440
per day at the testbed cadence and 48 per day at mainnet cadence, so the
accelerated testbed exaggerates the baseline flapping thirty-fold.

The model also predicts a limit of the current parameters at mainnet cadence.
When two Bitcoin blocks arrive within one polling interval, the next checkpoint
skips a block and the gap stays at 5 for a full Bitcoin interval, much longer
than \texttt{MaxSuspiciousTime} (24 Engram blocks, about 29\,s, chosen for the
20\,s cadence), so the FSM escalates to \textsc{Sovereign} without a fault. With
exponential block intervals a healthy chain does so about once in eight days
(10 events in 77 days of exposure, the same for $k{=}2$ and $k{=}4$) and never
with steady blocks or at the testbed cadence. This is a prediction of the
model, not an observation. The timeout must therefore follow the Bitcoin
cadence of the deployment: the live value of 24 blocks is 1.4 Bitcoin intervals
at the testbed's 20\,s, and a value of about 1.5 intervals (750 blocks, 15\,min,
at 600\,s) removes most false escalations at the price of about 14 more minutes
to reach \textsc{Sovereign} in a real fault (Table~\ref{tab:mst}). We keep 24 as the
default, which the testbed experiments use, and treat the timeout as a deployment
parameter. A live check at a 600\,s mining interval (steady blocks, three runs on fresh chains) reached \textsc{Sovereign} after 50.6, 50.9 and 60.7\,min (5, 6 and 6 Bitcoin blocks; 2{,}451, 2{,}477 and 2{,}921 Engram blocks), 54.1\,min on average. Two runs lie inside the range of the model (41 to 58\,min) and the third above it, close to the plain conversion of Table~\ref{tab:exposure}.

\begin{table}[h]
\caption{\texttt{MaxSuspiciousTime} at mainnet cadence (simulation; exponential
Bitcoin block intervals of mean 600\,s, $k{=}4$): false escalations of a healthy
chain per day (30 runs of up to 72\,h; the counts behind the small rates are 10,
3 and 1 events, so they indicate order of magnitude only) and the mean time from
a checkpoint pause to \textsc{Sovereign} (40 runs).}
\label{tab:mst}
\centering
\footnotesize
\begin{tabular}{@{}rrrr@{}}
\toprule
\textbf{Blocks} & \textbf{Time} & \textbf{False escalations per day} & \textbf{Mean to \textsc{Sovereign}, min} \\
\midrule
24 & 29\,s & 0.13 & 53 \\
500 & 10\,min & 0.04 & 63 \\
750 & 15\,min & 0.01 & 67 \\
1500 & 30\,min & 0.01 & 76 \\
\bottomrule
\end{tabular}
\end{table}

Table~\ref{tab:exposure-design} lists ways to bound the exposure and what each
costs. We implement none of them; the first two bound the number of withdrawals
without a lower threshold, and the last one shortens the window.

\begin{table}[h]
\caption{Ways to bound the exposure before the lock.}
\label{tab:exposure-design}
\centering
\footnotesize
\begin{tabular}{@{}p{0.25\linewidth}p{0.3\linewidth}p{0.3\linewidth}@{}}
\toprule
\textbf{Mechanism} & \textbf{Effect} & \textbf{Cost} \\
\midrule
Withdrawal quota per Bitcoin block, graded by the gap & $N \le$ (Bitcoin blocks in the window) $\times$ quota, independent of the Engram block rate and of the sender & Lower withdrawal throughput also when healthy \\
Release funds only after the block is anchored $K$ deep & Exposure close to zero & Every withdrawal waits about three Bitcoin blocks; needs a real bridge \\
Sensor on the age of the pending checkpoint against the expected cycle & Window shortens from about six to three or four Bitcoin blocks (our estimate) & Changes the FSM input, so every experiment on the Bitcoin gap would be rerun \\
\bottomrule
\end{tabular}
\end{table}


\subsection*{Hysteresis and concealment}

\textbf{Concealment against noise.} A higher down threshold $k$
tolerates longer bursts of noise, and it also tolerates a longer run of reports
that a Byzantine proposer interleaves with honest ones. Within the assumption
the choice of $k$ does not matter: for each of the 21 placements of two
Byzantine proposers in the seven-slot rotation, a real DA outage was detected
after exactly $k$ blocks ($k\in\{2,3,4\}$, simulation). With three Byzantine
validators, outside the assumption, a Byzantine proposer can keep a certificate
that reports a healthy DA layer, and an honest proposer then reports the outage
in the blocks it proposes; the outage is detected only if enough honest
proposers follow one another to reach $k$. In the simulation it was detected for all 35
placements at $k{=}2$, missed for 14 at $k{=}3$ and for 28 at $k{=}4$, and with
four Byzantine validators it was missed for every placement. The live check at
$k{=}2$ agrees: with three Byzantine validators the outage was detected in 5 of
5 runs. The adaptive threshold exposes the same
weakness only while its flap level is high (0 of 35 placements hidden at level 0,
28 at level 1, all 35 at level 2), whereas a static threshold of 4 or more is
exposed permanently.

\textbf{A window rule.} A rule that demotes when at least $m$ of the last $w$ blocks are bad,
instead of after a consecutive run, is not cut by interleaved reports. In a
reduced model of the \textsc{Anchored}$\leftrightarrow$\textsc{Suspicious} edge,
validated against the keeper-based simulation above (static thresholds and the
adaptive $2\to4$ reproduce its flaps, detection delays and hidden placements),
4 of 8 hid the outage in 0 of 35 placements, against 28 for a static threshold
of 4, at 0.8 flaps against 0.2 and the same 4-block detection. It does not
remove the cost of being fast: its detection is as slow as the static
threshold, and an adversary that sends bursts of $m$ blocks still obtains 200
demotions in 1200 blocks, where the adaptive threshold allows 1. Combining the
two (an adaptive window) gave 4.6 flaps and 172 demotions, so no single rule
won on all four criteria. We did not implement the window rule.

\textbf{Live check of concealment.} The live check at $k{=}4$ (three Byzantine validators with biased certificates, all seven bridges stopped for 120\,s, a fresh chain per run)
agrees with the simulation for adjacent placements and not for interleaved ones. With three consecutive Byzantine slots
the outage was detected in 3 of 3 runs, as predicted. With slots 0, 2 and 4 of the rotation, which the simulation predicts
would hide it, the chain also left \textsc{Anchored} and reached \textsc{Sovereign} in 3 of 3 runs, with 6 to 9 blocks outside the honest range
per run (6 to 12 over all six). We have not isolated why the prediction failed for the interleaved placement, so the concealment at $k{=}4$ rests on the simulation alone.

\textbf{Choosing the down threshold.} The experiments ran with
\texttt{DownHysteresisThreshold} of 2, and the live runs show that value to be too
low for a wide-area network: under the three-phase noise a static $k{=}2$ flapped
7.5 times and spent 82\% of the noisy phase in \textsc{Anchored}, a static $k{=}4$
did not flap and stayed \textsc{Anchored} throughout, at a detection delay of 13.7\,s
against 10.9\,s (an earlier single-run spot-check at $k{=}4$ did flap,
Section~\ref{subsec:rq3}). At the testbed's accelerated Bitcoin cadence the healthy checkpoint
cycle reaches the suspicious threshold in 3--16\% of the samples, and $k{=}2$ spent only
0.79--0.87 of the quiet phase in \textsc{Anchored}, against 1.00 for $k{=}4$; the
model of the cycle accounts for this, though we did not isolate it live.
We therefore recommend $k_{\downarrow} \geq 4$ as a deployment setting, set from the
noise of the environment like \texttt{MaxSuspiciousTime} is set from the Bitcoin
cadence, and keep 2 as the value of the reported experiments. The price is the
concealment of the preceding paragraph: an outage is detected only if $k$
consecutive proposers are honest, so the larger the threshold, the easier it is to
hide one.


\subsection*{Proving latency of the recovery proof}

\textbf{Measured proving latency.} In the first live deployment with chunk
aggregation, the prover built a 512-header aggregate from two chunks in about
4.3 minutes and a 514-header aggregate from three in about 6.6 minutes; the
tip then stood 211 and 321 blocks beyond the last covered header, close to the
product of latency and block rate. Both proofs were accepted by
\texttt{SubmitRecoveryProof}, but with the bound at 128 neither made the
recovery flag hold, and the chain stayed in \textsc{Recovering}. A single
256-header proof takes about two minutes, about 110 blocks, and fits, but
with little margin: at a block rate about 17\% higher, or a slower prover, it
would not, so the recovery times of the proof-only runs hold for this
testbed's block rate and proving speed. With the bound raised to 512 blocks and
everything else unchanged, the chain left \textsc{Recovering} through an
aggregate proof in both of two deployments. In the first the prover built a
514-header proof from three chunks in 6.5 minutes, and the tip stood 309
blocks (about 6.4 minutes) beyond the last covered header when it was
submitted; in the second a 512-header proof from two chunks took 4.6 minutes,
with the tip 224 blocks (about 4.7 minutes) beyond. All seven validators were
then in \textsc{Anchored}, 1{,}023 and 888 seconds after the chain started.
Proof-only recovery was on, so only the proof could have restored it. The
interval proved in these runs is the one since genesis, not one following an
outage injected into a healthy chain, and the runs with the bound at 128 are a
single deployment.

Aggregation bounds the verifier's cost, not the prover's. The aggregate is one
more proof on top of the chunk proofs, so building it takes longer than building the
chunks alone (49\% longer at 512 headers and 22\% at 1024 in
Appendix~\ref{app:aggregation}). The chunk proofs are independent given the chunk
boundary roots, so they can be built concurrently and joined by the aggregator
circuit.

\textbf{Proving latency on dedicated hosts.} Table~\ref{tab:instances} reports the proving step
(\texttt{nargo execute} and \texttt{bb prove} of one segment, in the published prover image) on three
\texttt{amd64} instance types, each alone on its host with no chain running. The witness generation is
single-threaded and takes 43\,s on the two newer types and 74\,s on the older one; \texttt{bb} uses
all eight cores. We chose the \texttt{c7a.2xlarge} for the proving workers of Testbed~B from this table, that is,
for its speed, before running the recovery experiments; the first run on each host, 4--7\,s slower (cold cache), is
left out of the means. On the multi-host chain the latency from the start of a proof to its accepted submission was
77\,s with one worker and 76\,s with two (73--87\,s per proof): the proving step as seen by the prover host,
69--75\,s including the transfer to the worker, and 2--17\,s of witness preparation, publication of the audit
trail and submission.

\begin{table}[h]
\caption{The proving step of one segment on one host, mean of the runs after the first (two runs on the
\texttt{c7a.2xlarge} and \texttt{c7i.2xlarge}, three on the \texttt{c5.2xlarge}; one host per type).}
\label{tab:instances}
\centering
\footnotesize
\begin{tabular}{@{}llrrr@{}}
\toprule
\textbf{Instance} & \textbf{CPU} & \textbf{\texttt{nargo execute} (s)} & \textbf{\texttt{bb prove} (s)} & \textbf{Step (s)} \\
\midrule
c7a.2xlarge & AMD EPYC 9R14 & 43.2 & 21.4 & 64.7 \\
c7i.2xlarge & Intel Xeon Platinum 8488C & 43.0 & 27.7 & 70.7 \\
c5.2xlarge & Intel Xeon Platinum 8124M & 74.0 & 37.1 & 111.0 \\
\bottomrule
\end{tabular}
\end{table}

Table~\ref{tab:parallel} reports the cost on the ten-core host with the testbed
stopped, at 0.8 blocks per second. Proving and aggregation alone fit inside the bound
of 128 blocks, with 65 blocks for 512 headers and 109 for 1024. The witnesses, which
are prepared one chunk after another, take 108\,s and 213\,s and are not part of that
figure. With them the latency from the end of the interval to an accepted proof is
about 189\,s (151 blocks) for 512 headers and 349\,s (280 blocks) for 1024, above the
bound. Two concurrent provers with five threads each shorten the proving step by 11\%
and 13\% because they share the cores; four concurrent provers did not start, since the
host lacks the memory. A single leaf proof takes 28\,s with all ten cores, 36\,s with
five and 60\,s with two.

\begin{table}[h]
\caption{Latency of building the recovery proof on one ten-core host (testbed
stopped). Proving is the chunk proofs plus the aggregation; the tail is the proving
time at 0.8 blocks per second. Sequential and parallel are means of two runs, except
the parallel run at 1024 headers, which is one run.}
\label{tab:parallel}
\centering
\footnotesize
\begin{tabular}{@{}rrrrrr@{}}
\toprule
 & & \multicolumn{2}{c}{\textbf{Sequential}} & \multicolumn{2}{c}{\textbf{Two provers}} \\
\cmidrule(lr){3-4}\cmidrule(lr){5-6}
\textbf{Headers} & \textbf{Witnesses (s)} & \textbf{Proving (s)} & \textbf{Tail} & \textbf{Proving (s)} & \textbf{Tail} \\
\midrule
512  & 108 & 81.3  & 65  & 72.6  & 58 \\
1024 & 213 & 136.5 & 109 & 118.9 & 95 \\
\bottomrule
\end{tabular}
\end{table}

Whether the prover keeps up with block production is therefore not established. It
needs a prover with cores of its own, and witnesses prepared chunk by chunk while the
chain is still in \textsc{Sovereign}, so that only the last chunk and the aggregation
remain on the critical path; we measured neither. The 4.3--6.6 minutes of the live
deployment were measured while the validators ran on the same cores, so the latencies
depend on the load of the host. The validators verify the submitted proof during
block execution, so the verification cost does not change with how the proof was
built.


\subsection*{Restart after a long chain}

The publisher's contiguous height is
kept in memory. A validator that restarts therefore republishes the chain from
height 0 and reports DA as failed until it has caught up. In one run, four
validators restarted on a chain of about 7{,}400 heights; the certified DA
health stayed failed and the chain remained in \textsc{Sovereign} for more than
half an hour. Persisting the contiguous height would remove the dependence on
chain length; the prototype does not.

